%% ============================================================
%%  Quantum Physics-Informed Neural Networks (QPINNs)
%%  Solving PDEs and ODEs with Variational Quantum Circuits
%%  Author: Morten Hjorth-Jensen
%%  Department of Physics, University of Oslo
%% ============================================================
\documentclass[aspectratio=169,10pt]{beamer}

%% -- Theme & colours ---------------------------------------------
\usetheme{Madrid}
\usecolortheme{seahorse}
\usefonttheme[onlymath]{serif}

\setbeamercolor{frametitle}{bg=blue!15!white,fg=blue!60!black}
\setbeamercolor{title}{fg=blue!70!black}
\setbeamercolor{block title}{bg=blue!20!white,fg=blue!60!black}
\setbeamercolor{block body}{bg=blue!5!white}
\setbeamercolor{alerted text}{fg=red!60!black}
\setbeamertemplate{navigation symbols}{}
\setbeamertemplate{footline}{%
  \leavevmode\hbox{%
    \begin{beamercolorbox}[wd=.25\paperwidth,ht=2.25ex,dp=1ex,center]{author in head/foot}%
      \usebeamerfont{author in head/foot}\insertshortauthor
    \end{beamercolorbox}%
    \begin{beamercolorbox}[wd=.50\paperwidth,ht=2.25ex,dp=1ex,center]{title in head/foot}%
      \usebeamerfont{title in head/foot}\insertshorttitle
    \end{beamercolorbox}%
    \begin{beamercolorbox}[wd=.25\paperwidth,ht=2.25ex,dp=1ex,right]{date in head/foot}%
      \usebeamerfont{date in head/foot}\insertshortdate{}\hspace*{1em}%
      \insertframenumber{}/\inserttotalframenumber\hspace*{1ex}
    \end{beamercolorbox}}%
  \vskip0pt}

%% -- Packages ----------------------------------------------------
\usepackage[utf8]{inputenc}
\usepackage{amsmath,amssymb,bm,physics}
\usepackage{braket}
\usepackage{graphicx}
\usepackage{booktabs}
\usepackage{xcolor}
\usepackage{tikz}
\usetikzlibrary{positioning,arrows.meta,shapes.geometric,
                decorations.pathreplacing,fit}
\usepackage{quantikz}
\usepackage{listings}
\usepackage{hyperref}

%% -- Python listing style -----------------------------------------
\lstset{
  language=Python,
  basicstyle=\ttfamily\scriptsize,
  keywordstyle=\color{blue!70!black}\bfseries,
  commentstyle=\color{gray},
  stringstyle=\color{red!60!black},
  breaklines=true,
  frame=single,
  backgroundcolor=\color{blue!3!white},
  xleftmargin=4pt, xrightmargin=4pt,
  numbers=left,
  numberstyle=\ttfamily\tiny\color{gray},
  numbersep=5pt
}

%% -- Metadata ----------------------------------------------------
\title[Quantum PINNs]{Quantum Physics-Informed Neural Networks}
\subtitle{Solving PDEs and ODEs with Variational Quantum Circuits}
\author[M.~Hjorth-Jensen]{Morten Hjorth-Jensen}
\institute[UiO]{Department of Physics \& CCSE\\University of Oslo}
\date{May 6, 2026}

%% ============================================================
\begin{document}
%% ============================================================

\begin{frame}
  \titlepage
\end{frame}

\begin{frame}{Outline}
  \tableofcontents[hidesubsections]
\end{frame}

%% ============================================================
\section{Recap: Quantum Neural Networks}
%% ============================================================

\begin{frame}{QNN as a Variational Quantum State}
  \begin{columns}[T]
    \column{0.52\textwidth}
      The central object of a \textbf{Quantum Neural Network}:
      \[
        \ket{\psi(x,\theta)} = U(\theta)\,U(x)\,\ket{0}^{\otimes n}
      \]
      \[
        f(x,\theta) = \bra{\psi(x,\theta)}\,O\,\ket{\psi(x,\theta)}
      \]
      \begin{itemize}
        \item $U(x)$: \textbf{feature map} --- encodes the input
        \item $U(\theta)$: \textbf{variational ansatz} --- trainable
        \item $O$: \textbf{observable} --- defines the output scalar
      \end{itemize}
      \begin{block}{Pauli ansatz}
        \[
          U(\theta)=e^{-iH(\theta)},\quad
          H(\theta)=\sum_\alpha\theta_\alpha P_\alpha
        \]
      \end{block}
    \column{0.44\textwidth}
      \begin{center}
      \begin{quantikz}[row sep=0.45cm, column sep=0.42cm]
        \lstick{$\ket{0}$}
          & \gate{R_x(x_1)} & \gate{R_y(\theta_1)} & \ctrl{1}
          & \gate{R_z(\theta_3)} & \meter{} \\
        \lstick{$\ket{0}$}
          & \gate{R_x(x_2)} & \gate{R_y(\theta_2)} & \targ{}
          & \gate{R_z(\theta_4)} & \meter{}
      \end{quantikz}
      \end{center}
      \vspace{0.2cm}
      \begin{alertblock}{Key tool: parameter-shift rule}
        \[
          \frac{\partial f}{\partial\theta}
          = \tfrac{1}{2}\bigl[f(\theta+\tfrac{\pi}{2})
                              -f(\theta-\tfrac{\pi}{2})\bigr]
        \]
        Exact gradient, 2 circuit evaluations per parameter.
      \end{alertblock}
  \end{columns}
\end{frame}

\begin{frame}{QFI, TDVP, and Natural Gradient}
  \begin{columns}[T]
    \column{0.52\textwidth}
      \textbf{Quantum Fisher Information (QFI):}
      \[
        F_{ij} = 4\,\mathrm{Re}\!\left(
          \langle\partial_i\psi|\partial_j\psi\rangle
          - \langle\partial_i\psi|\psi\rangle
            \langle\psi|\partial_j\psi\rangle
        \right)
      \]
      Riemannian metric on the variational manifold.

      \medskip
      \textbf{Time-Dependent Variational Principle (TDVP):}
      \[
        \sum_j F_{ij}\,\dot\theta_j = C_i,\quad
        C_i = \mathrm{Im}\langle\partial_i\psi|H|\psi\rangle
      \]
      \textbf{Natural gradient descent:}
      \[
        \dot\theta = F^{-1}\nabla_\theta\mathcal{L}
      \]
    \column{0.44\textwidth}
      \begin{block}{BCH expansion}
        \[
          U^\dagger O\,U = O + i[H,O]
            + \tfrac{i^2}{2}[H,[H,O]] + \cdots
        \]
        Depth $\leftrightarrow$ hierarchy of correlations.
      \end{block}
      \vspace{0.2cm}
      \begin{alertblock}{Barren plateaus}
        Deep random circuits:
        \[
          \mathrm{Var}(\nabla_\theta\mathcal{L})\sim 2^{-n}
        \]
        Mitigation: problem-inspired shallow ansätze ---
        exactly what QPINNs do.
      \end{alertblock}
  \end{columns}
\end{frame}

\begin{frame}{From Discriminative QNNs to Physics-Informed QNNs}
  \begin{columns}[T]
    \column{0.52\textwidth}
      \begin{block}{Classical discriminative QNN}
        \begin{itemize}
          \item Input: data pair $(x_k, y_k)$
          \item Loss: $\mathcal{L}=\sum_k(f(x_k;\theta)-y_k)^2$
          \item Learns a mapping from observations
          \item No physical constraint enforced
        \end{itemize}
      \end{block}
      \vspace{0.3cm}
      \begin{alertblock}{Physics-informed QNN (QPINN)}
        \begin{itemize}
          \item Input: collocation points $x_k\in\Omega$
          \item Loss: residual of a differential equation
            \[
              \mathcal{L}=\|\mathcal{D}[f_\theta]\|^2
                         +\lambda\|\mathcal{B}[f_\theta]\|^2
            \]
          \item Learns the solution to a PDE / ODE
          \item Physics is baked into the loss function
        \end{itemize}
      \end{alertblock}
    \column{0.44\textwidth}
      \textbf{Why quantum circuits for PDEs?}
      \begin{enumerate}\itemsep4pt
        \item Potential exponential compression of solution
              space via entanglement
        \item Natural encoding of wave-like solutions
              through rotation gates
        \item Gradient information from parameter-shift rule
              maps directly onto PDE residuals
        \item Hardware noise acts as a form of regularisation
        \item Connection to quantum simulation of
              differential operators
      \end{enumerate}
  \end{columns}
\end{frame}

%% ============================================================
\section{Classical Physics-Informed Neural Networks}
%% ============================================================

\begin{frame}{Classical PINNs: Raissi, Perdikaris, Karniadakis (2019)}
  \begin{columns}[T]
    \column{0.52\textwidth}
      A classical PINN approximates the solution $u(x,t)$ of a PDE
      \[
        \mathcal{N}[u](x,t) = 0,\quad (x,t)\in\Omega
      \]
      by a neural network $u_\theta(x,t)$ trained to minimise:
      \[
        \mathcal{L}(\theta)
        = \underbrace{\frac{1}{N_r}\sum_{k=1}^{N_r}
            \bigl|\mathcal{N}[u_\theta](x_k,t_k)\bigr|^2}_{
            \text{PDE residual loss}\;\mathcal{L}_r}
          + \lambda\,\underbrace{\frac{1}{N_b}\sum_{k=1}^{N_b}
            \bigl|u_\theta(x_k^b,t_k^b)-g_k\bigr|^2}_{
            \text{BC/IC loss}\;\mathcal{L}_b}
      \]
      where $\{x_k\}$ are \emph{collocation points} sampled in $\Omega$
      and $\{x_k^b\}$ are boundary/initial condition points.
    \column{0.44\textwidth}
      \begin{block}{Automatic differentiation}
        Classical PINNs exploit automatic differentiation
        to evaluate $\mathcal{N}[u_\theta]$ exactly
        (no finite differences).
      \end{block}
      \vspace{0.2cm}
      \begin{block}{Examples}
        \begin{itemize}
          \item Poisson: $-\nabla^2 u = f$
          \item Heat: $\partial_t u - \kappa\nabla^2 u = 0$
          \item Wave: $\partial_{tt}u - c^2\nabla^2 u = 0$
          \item Burgers: $\partial_t u + u\partial_x u = \nu\partial_{xx}u$
          \item Schrödinger: $i\partial_t\psi = H\psi$
        \end{itemize}
      \end{block}
      \begin{alertblock}{Key challenge}
        Spectral bias: classical NNs learn low-frequency
        components of $u$ much faster than high-frequency ones.
      \end{alertblock}
  \end{columns}
\end{frame}

\begin{frame}{PINN Loss Landscape and Gradient Flow}
  \begin{columns}[T]
    \column{0.52\textwidth}
      For a general differential operator $\mathcal{D}$ the
      residual at collocation point $x_k$ is:
      \[
        r_k(\theta) = \mathcal{D}[f_\theta](x_k)
      \]
      The total loss and its gradient:
      \[
        \mathcal{L}(\theta)
        = \frac{1}{N}\sum_k r_k^2(\theta),\quad
        \nabla_\theta\mathcal{L}
        = \frac{2}{N}\sum_k r_k\,\nabla_\theta r_k
      \]
      The Neural Tangent Kernel (NTK) of the PINN governs
      training dynamics:
      \[
        \Theta_{kl}(\theta)
        = \nabla_\theta f_\theta(x_k)^\top
          \nabla_\theta f_\theta(x_l)
      \]
      Eigenvalues of $\Theta$ determine the convergence rate
      for each frequency component.
    \column{0.44\textwidth}
      \begin{alertblock}{Failure modes of classical PINNs}
        \begin{enumerate}
          \item \textbf{Spectral bias:} slow convergence for
                high-frequency solutions
          \item \textbf{Ill-conditioning:} loss balance
                between $\mathcal{L}_r$ and $\mathcal{L}_b$
          \item \textbf{Stiff equations:} chaotic loss landscape
          \item \textbf{High dimensionality:} curse of
                dimensionality for high-$d$ PDEs
        \end{enumerate}
      \end{alertblock}
      \vspace{0.2cm}
      \begin{block}{Quantum motivation}
        QPINNs address points 1 and 4: entangled circuits
        have richer frequency spectra; quantum memory scales
        exponentially with qubits.
      \end{block}
  \end{columns}
\end{frame}

%% ============================================================
\section{Quantum PINNs: Mathematical Foundation}
%% ============================================================

\begin{frame}{QPINN: The Quantum Ansatz for PDE Solutions}
  \begin{columns}[T]
    \column{0.52\textwidth}
      Replace the classical network $u_\theta(x)$ with a
      \textbf{variational quantum circuit output}:
      \[
        u_\theta(x) = \bra{0}U^\dagger(x,\theta)\,O\,U(x,\theta)\ket{0}
      \]
      \[
        = \bra{\psi(x,\theta)}\,O\,\ket{\psi(x,\theta)}
      \]
      where:
      \begin{itemize}
        \item $U(x,\theta) = W_L(\theta)\cdots W_1(\theta)\,S(x)$
        \item $S(x)$: data-encoding unitary (feature map)
        \item $W_l(\theta)$: variational layers
        \item $O$: a fixed observable (e.g., $Z_0$)
      \end{itemize}
      The output $u_\theta(x)\in[-1,1]$ approximates the
      normalised PDE solution.
    \column{0.44\textwidth}
      \begin{center}
      \begin{quantikz}[row sep=0.42cm, column sep=0.38cm]
        \lstick{$\ket{0}$}
          & \gate{S_1(x)} & \gate{W_1(\theta)} & \ctrl{1}
          & \gate{W_1'(\theta)} & \meter{$O$} \\
        \lstick{$\ket{0}$}
          & \gate{S_2(x)} & \gate{W_2(\theta)} & \targ{}
          & \gate{W_2'(\theta)} & \qw \\
        \lstick{$\ket{0}$}
          & \gate{S_3(x)} & \gate{W_3(\theta)} & \ctrl{1}
          & \gate{W_3'(\theta)} & \qw \\
        \lstick{$\ket{0}$}
          & \gate{S_4(x)} & \gate{W_4(\theta)} & \targ{}
          & \gate{W_4'(\theta)} & \qw
      \end{quantikz}
      \end{center}
      {\small $S_i(x)$: encoding of coordinate $x$;\\
       $W_i(\theta)$: variational Pauli rotations;\\
       $O=Z_0$: observable giving scalar output.}
  \end{columns}
\end{frame}

\begin{frame}{Data Encoding Strategies for PDEs}
  The choice of encoding $x\mapsto S(x)$ is critical for QPINNs.

  \begin{columns}[T]
    \column{0.48\textwidth}
      \textbf{1.\ Angle encoding}
      \[
        S(x) = \bigotimes_{i=1}^n R_y(x\pi)
      \]
      Simple, one qubit per input; output is a
      degree-$n$ Fourier series in $x$.

      \vspace{0.3cm}
      \textbf{2.\ Repeated angle encoding}
      \[
        S(x) = \prod_{l=1}^L\bigl[W_l(\theta)\,
          {\textstyle\bigotimes_i} R_y(x)\bigr]
      \]
      Interleaving data re-uploads with variational layers
      (\emph{data re-uploading}): extends the expressible
      frequency spectrum to $[-Ln/2, Ln/2]$.
    \column{0.48\textwidth}
      \textbf{3.\ IQP / ZZ encoding} (richer correlations)
      \[
        S(x) = e^{-ixH_{\rm enc}},\quad
        H_{\rm enc} = \sum_i x\,Z_i + \sum_{i<j}x^2 Z_iZ_j
      \]

      \vspace{0.3cm}
      \begin{alertblock}{Fourier perspective (Schuld et al.)}
        Any QNN output is a partial Fourier series:
        \[
          f_\theta(x) = \sum_{\omega\in\Omega}
            c_\omega(\theta)\,e^{i\omega x}
        \]
        The set of accessible frequencies $\Omega$ is
        determined entirely by the data-encoding gates.
        Wider $\Omega$ $\Rightarrow$ richer PDE solutions.
      \end{alertblock}
  \end{columns}
\end{frame}

\begin{frame}{Derivatives of the Quantum Ansatz}
  \textbf{The core challenge:} PDE residuals require
  $\partial_x u_\theta$, $\partial_{xx} u_\theta$, etc.
  These are \emph{derivatives with respect to the input $x$},
  not the parameters $\theta$.

  \begin{columns}[T]
    \column{0.52\textwidth}
      \textbf{Input-shift rule.}
      For angle encoding $S(x) = e^{-ixP/2}$ (Pauli $P$):
      \[
        \frac{\partial u_\theta}{\partial x}
        = \tfrac{1}{2}\bigl[
            u_\theta(x+\tfrac{\pi}{2}) - u_\theta(x-\tfrac{\pi}{2})
          \bigr]
      \]
      This is the \emph{same} parameter-shift formula, now
      applied to the \textbf{input encoding} gate.

      \textbf{Second derivative:}
      \begin{align*}
        \frac{\partial^2 u_\theta}{\partial x^2}
        &= \tfrac{1}{2}\bigl[
            u_\theta(x+\pi) - 2u_\theta(x) + u_\theta(x-\pi)
          \bigr]
      \end{align*}
      General $n$-th order derivative via repeated shifts.
    \column{0.44\textwidth}
      \begin{alertblock}{Key result}
        All spatial derivatives of $u_\theta(x)$ are
        computable by evaluating the \emph{same circuit} at
        shifted input values --- no ancilla qubits, no
        additional circuit depth.
      \end{alertblock}
      \vspace{0.2cm}
      \begin{block}{Multi-dimensional PDE}
        For $u(x_1,\ldots,x_d)$ with $n_i$ qubits encoding
        $x_i$:
        \[
          \partial_{x_i}\partial_{x_j}u_\theta
          = \text{shift in } x_i \text{ and } x_j
        \]
        Cross-derivatives come from shifts in
        \emph{two} inputs simultaneously.
      \end{block}
  \end{columns}
\end{frame}

\begin{frame}{The QPINN Loss Function}
  \begin{columns}[T]
    \column{0.52\textwidth}
      For a PDE $\mathcal{D}[u](x) = f(x)$ on $\Omega$
      with boundary conditions $u(x) = g(x)$ on $\partial\Omega$:

      \textbf{Total QPINN loss:}
      \[
        \mathcal{L}(\theta)
        = \mathcal{L}_r(\theta) + \lambda\,\mathcal{L}_b(\theta)
      \]
      \[
        \mathcal{L}_r = \frac{1}{N_r}\sum_{k=1}^{N_r}
          \bigl|\mathcal{D}[u_\theta](x_k) - f(x_k)\bigr|^2
      \]
      \[
        \mathcal{L}_b = \frac{1}{N_b}\sum_{k=1}^{N_b}
          \bigl|u_\theta(x_k^b) - g(x_k^b)\bigr|^2
      \]
      where $u_\theta(x) = \bra{\psi(x,\theta)}O\ket{\psi(x,\theta)}$
      and all derivatives entering $\mathcal{D}[u_\theta]$ are
      evaluated via input-shift rules.
    \column{0.44\textwidth}
      \begin{block}{Gradient w.r.t.\ $\theta_\mu$}
        \[
          \frac{\partial\mathcal{L}_r}{\partial\theta_\mu}
          = \frac{2}{N_r}\sum_k r_k\,
            \frac{\partial r_k}{\partial\theta_\mu}
        \]
        where $r_k = \mathcal{D}[u_\theta](x_k) - f(x_k)$
        and
        \[
          \frac{\partial r_k}{\partial\theta_\mu}
          = \mathcal{D}\!\left[
              \frac{\partial u_\theta}{\partial\theta_\mu}
            \right](x_k)
        \]
        The operator $\mathcal{D}$ is linear in $u_\theta$
        $\Rightarrow$ shifts and $\mathcal{D}$ commute.
      \end{block}
      \begin{alertblock}{Parameter-shift + input-shift}
        Both types of derivatives are computable with
        $O(N_p\cdot N_r)$ circuit evaluations.
      \end{alertblock}
  \end{columns}
\end{frame}

\begin{frame}{QPINN Circuit for a 1D PDE}
  Complete QPINN circuit for $-\partial_{xx}u = f(x)$,
  $x\in[0,1]$, with $n=4$ qubits:

  \vspace{0.2cm}
  \begin{center}
  \begin{quantikz}[row sep=0.42cm, column sep=0.38cm]
    \lstick{$\ket{0}$}
      & \gate{R_y(\pi x)} & \gate{R_y(\theta_1^{(1)})} & \ctrl{1} & \qw
      & \gate{R_y(\theta_1^{(2)})} & \ctrl{1} & \qw
      & \meter{$Z$} \\
    \lstick{$\ket{0}$}
      & \gate{R_y(\pi x)} & \gate{R_y(\theta_2^{(1)})} & \targ{} & \ctrl{1}
      & \gate{R_y(\theta_2^{(2)})} & \targ{} & \ctrl{1}
      & \qw \\
    \lstick{$\ket{0}$}
      & \gate{R_y(\pi x)} & \gate{R_y(\theta_3^{(1)})} & \qw & \targ{}
      & \gate{R_y(\theta_3^{(2)})} & \qw & \targ{}
      & \qw \\
    \lstick{$\ket{0}$}
      & \gate{R_y(\pi x)} & \gate{R_y(\theta_4^{(1)})} & \qw & \qw
      & \gate{R_y(\theta_4^{(2)})} & \qw & \qw
      & \qw
  \end{quantikz}
  \end{center}

  \vspace{0.2cm}
  \begin{columns}[T]
    \column{0.48\textwidth}
      \textbf{Evaluation for residual:}
      \begin{enumerate}
        \item $u_\theta(x)$: run circuit with input $x$
        \item $u_\theta(x\pm\frac\pi2)$: two shifted circuits
              for $\partial_{xx}u_\theta$ via second-order shift
      \end{enumerate}
    \column{0.48\textwidth}
      \textbf{Gradient for training:}
      \begin{enumerate}
        \item For each $\theta_\mu$: parameter-shift at
              $x$, $x\pm\frac\pi2$
        \item Total: $2N_p \times 3N_r$ circuit evaluations
              per gradient step
      \end{enumerate}
  \end{columns}
\end{frame}

%% ============================================================
\section{Poisson Equation and Quantum Solution}
%% ============================================================

\begin{frame}{The 1D Poisson Equation}
  \begin{columns}[T]
    \column{0.52\textwidth}
      \textbf{Problem statement:}
      \[
        -\frac{d^2 u}{dx^2} = f(x),\quad x\in(0,1)
      \]
      \[
        u(0) = u(1) = 0 \quad\text{(Dirichlet BCs)}
      \]
      \textbf{Exact solution} for $f(x) = \pi^2\sin(\pi x)$:
      \[
        u^\star(x) = \sin(\pi x)
      \]

      \textbf{QPINN residual loss:}
      \[
        \mathcal{L}_r = \frac{1}{N_r}\sum_{k=1}^{N_r}
          \left[-\frac{d^2 u_\theta}{dx^2}\bigg|_{x_k}
                - f(x_k)\right]^2
      \]
      \textbf{Boundary loss:}
      \[
        \mathcal{L}_b = u_\theta(0)^2 + u_\theta(1)^2
      \]
    \column{0.44\textwidth}
      \textbf{Second derivative via input-shift:}
      \[
        \frac{d^2 u_\theta}{dx^2}\bigg|_x
        = \frac{1}{2}\left[
            u_\theta(x{+}\tfrac\pi2)
            - 2u_\theta(x)
            + u_\theta(x{-}\tfrac\pi2)
          \right]
      \]
      This uses the identity for $u_\theta(x)
      = \langle O\rangle_{x}$
      with $S(x) = e^{-ixP/2}$:
      \[
        \partial_{xx}\langle O\rangle_x
        = -\langle O\rangle_x + \tfrac12\bigl(
          \langle O\rangle_{x+\pi/2}
          + \langle O\rangle_{x-\pi/2}\bigr)
      \]
      \begin{alertblock}{Three circuit evaluations}
        The Laplacian at one point costs exactly 3 circuit
        runs: $x-\pi/2$, $x$, $x+\pi/2$.
      \end{alertblock}
  \end{columns}
\end{frame}

\begin{frame}{Variational Formulation and Connection to FEM}
  \begin{columns}[T]
    \column{0.52\textwidth}
      The Poisson equation has an equivalent \textbf{variational form}:
      \[
        \min_u\,\mathcal{E}[u]
        = \int_0^1\!\left[\tfrac12(u')^2 - f\,u\right]dx
      \]
      subject to $u(0)=u(1)=0$.

      The QPINN can be trained directly on the
      \emph{energy functional} (Ritz method):
      \[
        \mathcal{L}_{\rm Ritz}(\theta)
        = \sum_k w_k\!\left[
            \tfrac12\left(\frac{du_\theta}{dx}\bigg|_{x_k}\right)^2
            - f(x_k)\,u_\theta(x_k)
          \right]
      \]
      where $w_k$ are quadrature weights.
    \column{0.44\textwidth}
      \begin{block}{Why Ritz can outperform strong form}
        \begin{itemize}
          \item Needs only first derivatives
                (one input-shift per point)
          \item Smoother loss landscape
          \item Natural for symmetric positive-definite operators
          \item Directly minimises the physical energy
        \end{itemize}
      \end{block}
      \vspace{0.2cm}
      \begin{alertblock}{Connection to FEM}
        The Ritz-QPINN is a \emph{quantum Galerkin method}:
        the quantum ansatz $u_\theta(x)$ plays the role of
        the finite element basis, but spans an exponentially
        large function space.
      \end{alertblock}
  \end{columns}
\end{frame}

\begin{frame}{Generalisation: ODEs and Other PDEs}
  \begin{columns}[T]
    \column{0.52\textwidth}
      \textbf{1D harmonic oscillator ODE:}
      \[
        \frac{d^2u}{dx^2} + \omega^2 u = 0
      \]
      Residual: $r_k = u_\theta''(x_k) + \omega^2 u_\theta(x_k)$

      \medskip
      \textbf{1D heat equation (time-dependent):}
      \[
        \partial_t u - \kappa\partial_{xx}u = 0
      \]
      Encode $(x,t)$ into circuit:
      $S(x,t)=R_y(\pi x)\otimes R_y(\pi t)$;
      evaluate time and space derivatives by input-shift
      on respective qubit registers.

      \medskip
      \textbf{Burgers equation (nonlinear):}
      \[
        \partial_t u + u\,\partial_x u = \nu\,\partial_{xx}u
      \]
      Nonlinear term $u\partial_x u$: product of two
      circuit outputs at shifted inputs --- still computable.
    \column{0.44\textwidth}
      \begin{block}{Schrödinger equation as a QPINN}
        \[
          i\partial_t\psi = -\tfrac{1}{2}\partial_{xx}\psi
                            + V(x)\psi
        \]
        The quantum ansatz $\psi_\theta(x,t)$ has a
        \emph{natural} complex structure:
        measure real and imaginary parts via two
        different observables ($Z$ and $Y$).

        The QPINN loss enforces the time-evolution equation
        directly --- a bridge between QPINNs and
        Hamiltonian simulation.
      \end{block}
      \begin{alertblock}{Key point}
        All linear differential operators act linearly
        on the circuit output $u_\theta(x)$, so residuals
        always reduce to sums of circuit evaluations at
        shifted inputs.
      \end{alertblock}
  \end{columns}
\end{frame}

%% ============================================================
\section{Derivative Computation via Parameter and Input Shifts}
%% ============================================================

\begin{frame}{Unified Shift Framework: Parameters and Inputs}
  \begin{columns}[T]
    \column{0.52\textwidth}
      For a quantum circuit with a gate $e^{-i\phi P/2}$
      (Pauli $P$, half-angle $\phi$):
      \[
        \frac{\partial}{\partial\phi}
        \langle O\rangle_\phi
        = \frac{1}{2}\bigl[
            \langle O\rangle_{\phi+\pi/2}
            - \langle O\rangle_{\phi-\pi/2}
          \bigr]
      \]
      \textbf{This formula is universal:}
      \begin{itemize}
        \item $\phi = \theta_\mu$ (variational parameter):
              \textbf{parameter-shift rule} --- gives gradient
              w.r.t.\ trainable weights
        \item $\phi = x$ (input encoding):
              \textbf{input-shift rule} --- gives spatial
              derivative of the PDE solution
      \end{itemize}

      Higher-order derivatives by repeated application:
      \[
        \partial^n_\phi\langle O\rangle
        = \sum_{k=0}^n(-1)^k\binom{n}{k}
          \langle O\rangle_{\phi+(\frac{n}{2}-k)\frac\pi2}
      \]
    \column{0.44\textwidth}
      \begin{center}
      \begin{quantikz}[row sep=0.45cm, column sep=0.38cm]
        \lstick{$\ket{0}$}
          & \gate{R_y(x\pm\frac\pi2)} & \gate{W(\theta)} & \meter{}
      \end{quantikz}
      \end{center}
      {\small Input-shift: evaluates $\partial_x u_\theta$
       using $3$ circuit runs at $x{-}\frac\pi2$, $x$, $x{+}\frac\pi2$.}

      \vspace{0.4cm}
      \begin{center}
      \begin{quantikz}[row sep=0.45cm, column sep=0.38cm]
        \lstick{$\ket{0}$}
          & \gate{R_y(x)} & \gate{R_y(\theta\pm\frac\pi2)} & \meter{}
      \end{quantikz}
      \end{center}
      {\small Parameter-shift: evaluates $\partial_\theta u_\theta$
       using $2$ circuit runs at $\theta{\pm}\frac\pi2$.}

      \begin{alertblock}{Total budget}
        $k$-th order spatial derivative at one point:
        $k+1$ circuit runs.
      \end{alertblock}
  \end{columns}
\end{frame}

\begin{frame}{Commutation of Differential Operators with Circuit Output}
  \begin{columns}[T]
    \column{0.52\textwidth}
      Let $u_\theta(x) = \bra{0}U^\dagger(x,\theta)O\,U(x,\theta)\ket{0}$.
      For a \emph{linear} differential operator $\mathcal{D}$:
      \[
        \mathcal{D}[u_\theta](x)
        = \bra{0}U^\dagger(x,\theta)\,
          \tilde{O}(x,\mathcal{D})\,
          U(x,\theta)\ket{0}
      \]
      where $\tilde{O}(x,\mathcal{D})$ is the \emph{effective
      observable} induced by $\mathcal{D}$ acting on $U(x,\theta)$.

      \textbf{For $\mathcal{D} = \partial_x^2$:}
      \[
        \partial_{xx} u_\theta(x)
        = \bra{0}\partial_{xx}\!\left[
            U^\dagger O\,U
          \right]\ket{0}
        = \bra{0}\tilde{O}_{xx}\ket{0}
      \]
      The effective observable $\tilde{O}_{xx}$ is:
      \[
        \tilde{O}_{xx}
        = U^\dagger\!(\partial_{xx} U)O
          + 2U^\dagger\!(\partial_x U)O(\partial_x U^\dagger)U
          + \cdots
      \]
    \column{0.44\textwidth}
      \begin{block}{Practical evaluation}
        Rather than computing $\tilde{O}_{xx}$ analytically,
        the input-shift rule evaluates it implicitly:
        \[
          \partial_{xx} u_\theta(x)
          \approx \frac{u_\theta(x+h)-2u_\theta(x)+u_\theta(x-h)}{h^2}
        \]
        with $h=\pi/2$ for \emph{exact} (not approximate) results.
      \end{block}
      \vspace{0.2cm}
      \begin{alertblock}{Linearity is key}
        For \emph{linear} PDEs, the QPINN gradient decomposes
        into a sum of parameter-shift evaluations, each at a
        different input shift --- no cross-terms.
        Nonlinear PDEs require products of circuit outputs.
      \end{alertblock}
  \end{columns}
\end{frame}

%% ============================================================
\section{Expressivity, Error Bounds, and Barren Plateaus}
%% ============================================================

\begin{frame}{Expressivity of QPINNs: Fourier Analysis}
  \begin{columns}[T]
    \column{0.52\textwidth}
      \textbf{Theorem (Schuld et al.\ 2021).}
      The output of any QNN with data-encoding gates
      $e^{-ixP/2}$ (Pauli $P$) is a \emph{partial Fourier series}:
      \[
        u_\theta(x) = \sum_{\omega\in\Omega}
          c_\omega(\theta)\,e^{i\omega x}
      \]
      where the set of accessible frequencies is:
      \[
        \Omega \subseteq \bigl\{
          \omega_1 + \omega_2 + \cdots + \omega_n
          : \omega_i\in\{0,\pm\tfrac{1}{2}\}
        \bigr\}
      \]
      The coefficients $c_\omega(\theta)$ depend on the
      variational parameters; the frequencies $\Omega$
      depend only on the \emph{encoding}.

      \textbf{For solving PDEs,} the highest frequency
      in $\Omega$ limits the shortest spatial scale
      that can be resolved.
    \column{0.44\textwidth}
      \begin{block}{Extending the frequency spectrum}
        \begin{itemize}
          \item More qubits $\to$ denser $\Omega$
          \item Data re-uploading $\to$ integer multiples
          \item IQP encoding $\to$ pairwise frequencies
                $x_ix_j$
        \end{itemize}
      \end{block}
      \vspace{0.2cm}
      \begin{alertblock}{Consequence for PDEs}
        A QPINN with $n$ qubits and $L$ re-uploading layers
        can represent solutions containing frequencies up
        to $\omega_{\max} = Ln/2$.
        For a target function with frequencies up to
        $\omega^\star$, we need:
        \[
          Ln \geq 2\omega^\star
        \]
        qubits $\times$ layers.
      \end{alertblock}
  \end{columns}
\end{frame}

\begin{frame}{Approximation Error and Convergence}
  \begin{columns}[T]
    \column{0.52\textwidth}
      Let $u^\star\in H^s(\Omega)$ (Sobolev space of order $s$)
      be the exact PDE solution. The QPINN approximation error
      decomposes into three parts:
      \[
        \|u_{\theta^\star} - u^\star\|_{L^2}
        \leq \underbrace{\varepsilon_{\rm approx}}_{\text{expressivity}}
           + \underbrace{\varepsilon_{\rm opt}}_{\text{optimisation}}
           + \underbrace{\varepsilon_{\rm stat}}_{\text{sampling}}
      \]

      \textbf{Approximation error} (best possible):
      \[
        \varepsilon_{\rm approx}
        \sim \|\hat{u}^\star_\omega\|
             \text{ for }\omega\notin\Omega
      \]
      Controlled by the frequency spectrum of $u^\star$ and $\Omega$.

      \textbf{Optimisation error:} depends on the loss landscape
      (barren plateaus, local minima).

      \textbf{Statistical error:} $\varepsilon_{\rm stat}
      \sim N_r^{-1/2}$ from finite collocation points.
    \column{0.44\textwidth}
      \begin{block}{No free lunch}
        QPINNs do not automatically outperform classical
        PINNs for arbitrary PDEs.
        Advantage requires:
        \begin{enumerate}
          \item Solutions with quantum-compressible
                frequency structure
          \item High-dimensional problems ($d\gg 1$)
                where quantum memory helps
          \item Availability of quantum hardware
        \end{enumerate}
      \end{block}
      \vspace{0.2cm}
      \begin{alertblock}{High-$d$ advantage}
        For $d$-dimensional PDEs, the classical PINN needs
        $O(\varepsilon^{-d})$ parameters; the QPINN can
        compress the solution using $O(d)$ qubits if the
        solution has a tensor-product structure.
      \end{alertblock}
  \end{columns}
\end{frame}

\begin{frame}{Barren Plateaus in QPINNs and Mitigation}
  \begin{columns}[T]
    \column{0.52\textwidth}
      QPINNs suffer from the same barren plateau phenomenon
      as general QNNs (from week15.tex):
      \[
        \mathrm{Var}\!\left(
          \frac{\partial\mathcal{L}}{\partial\theta_\mu}
        \right)
        \sim \frac{1}{2^n}
      \]
      \textbf{Additional sources in QPINNs:}
      \begin{itemize}
        \item Residual $r_k\approx 0$ early in training:
              gradient is small even without barren plateaus
        \item Loss balance: $\mathcal{L}_r\gg\mathcal{L}_b$
              makes BC enforcement slow
        \item Stiff PDEs: widely varying eigenvalues of the
              differential operator
      \end{itemize}
    \column{0.44\textwidth}
      \begin{block}{Mitigation strategies for QPINNs}
        \begin{enumerate}
          \item \textbf{Physics-inspired ansatz:}
                layer structure mirrors the Green's function
                of the differential operator
          \item \textbf{Incremental training:}
                start with low-frequency modes,
                progressively add higher frequencies
          \item \textbf{Adaptive collocation:}
                concentrate points where residual is large
          \item \textbf{Loss weighting:}
                normalise $\mathcal{L}_r$ and $\mathcal{L}_b$
                by their gradient norms
          \item \textbf{Natural gradient (QFI):}
                \[
                  \dot\theta = F^{-1}\nabla_\theta\mathcal{L}
                \]
                Respects the geometry of the quantum
                state manifold
        \end{enumerate}
      \end{block}
  \end{columns}
\end{frame}

\begin{frame}{QPINN vs.\ Classical PINN: Comparison}
  \begin{center}
  \renewcommand{\arraystretch}{1.3}
  \begin{tabular}{p{3.2cm}p{4.8cm}p{4.8cm}}
    \toprule
    \textbf{Feature} & \textbf{Classical PINN} & \textbf{QPINN} \\
    \midrule
    Ansatz       & Deep MLP $u_\theta(x)$ &
                   Circuit output $\langle O\rangle_{x,\theta}$ \\
    Derivatives  & Automatic differentiation &
                   Input-shift rule (exact) \\
    Gradient     & Backpropagation &
                   Parameter-shift rule (exact) \\
    Memory       & $O(N_p)$ parameters &
                   $2^n$ amplitudes, $O(n)$ params \\
    Freq.\ spectrum & Arbitrary (but spectral bias) &
                   Bounded by encoding \\
    High-$d$ PDE & Exponential cost &
                   Potential $O(d)$ qubit advantage \\
    Hardware     & GPU/CPU & Quantum device or simulator \\
    Maturity     & Production-ready & Research / NISQ \\
    \bottomrule
  \end{tabular}
  \end{center}
  \vspace{0.2cm}
  \begin{alertblock}{When QPINNs are most promising}
    High-dimensional PDEs; solutions with structured
    frequency spectra; problems with natural quantum
    simulation analogues (e.g., Schrödinger, Maxwell).
  \end{alertblock}
\end{frame}

%% ============================================================
\section{Python Implementation: 1D Poisson Equation}
%% ============================================================

\begin{frame}[fragile]{Python Implementation: Gates and Circuit}
  \begin{lstlisting}
import numpy as np
from scipy.optimize import minimize

# -- Gate definitions --------------------------------------------
def ry(t):
    c, s = np.cos(t / 2), np.sin(t / 2)
    return np.array([[c, -s], [s, c]], dtype=complex)

def rz(t):
    return np.diag([np.exp(-1j * t / 2), np.exp(1j * t / 2)]).astype(complex)

I2   = np.eye(2, dtype=complex)
CNOT = np.array([[1,0,0,0],[0,1,0,0],[0,0,0,1],[0,0,1,0]], dtype=complex)
Z0   = np.kron(np.diag([1., -1.]).astype(complex), I2)  # Z on qubit 0

def var_layer(p):
    """Ry(p0)xRy(p1) -> CNOT -> Rz(p2)xRz(p3)"""
    return np.kron(rz(p[2]), rz(p[3])) @ CNOT @ np.kron(ry(p[0]), ry(p[1]))

def circuit(x, params):
    """
    2-qubit QPINN circuit for 1D PDE.
    Encoding: Ry(pi*x) on each qubit.
    Ansatz:   2 variational layers (8 parameters).
    Output:   <Z_0> in [-1, 1].
    """
    state = np.zeros(4, dtype=complex); state[0] = 1.0
    # Angle encoding: x -> Ry(pi*x)
    U = np.kron(ry(np.pi * x), ry(np.pi * x))
    U = var_layer(params[:4]) @ U    # variational layer 1
    U = var_layer(params[4:8]) @ U   # variational layer 2
    psi = U @ state
    return float((psi.conj() @ Z0 @ psi).real)
  \end{lstlisting}
\end{frame}

\begin{frame}[fragile]{Derivatives via Input-Shift Rule}
  \begin{lstlisting}
# -- Input-shift rule for spatial derivatives ---------------------
SHIFT = np.pi / 2   # shift for Pauli Ry generator

def du_dx(x, params):
    """First derivative d/dx u_theta(x) via input-shift rule."""
    return 0.5 * (circuit(x + SHIFT, params)
                - circuit(x - SHIFT, params))

def d2u_dx2(x, params):
    """
    Second derivative d^2/dx^2 u_theta(x) via second-order shift.
    Uses: f''(x) = f(x + pi/2) - 2*f(x) + f(x - pi/2)
    (exact for generators with eigenvalues +/- 1/2).
    """
    return (circuit(x + SHIFT, params)
          - 2 * circuit(x, params)
          + circuit(x - SHIFT, params))

# -- Source term for Poisson: f(x) = pi^2 * sin(pi*x) ------------
def source(x):
    return np.pi**2 * np.sin(np.pi * x)

# -- Exact solution for validation -------------------------------
def exact(x):
    return np.sin(np.pi * x)
  \end{lstlisting}
\end{frame}

\begin{frame}[fragile]{QPINN Loss Function and Gradient}
  \begin{lstlisting}
# -- Collocation and boundary points -----------------------------
np.random.seed(42)
N_r = 20    # interior collocation points
N_b = 2     # boundary points (x=0 and x=1)
x_r = np.linspace(0.05, 0.95, N_r)    # interior
x_b = np.array([0.0, 1.0])            # boundary
lambda_bc = 10.0                       # BC weight

# -- QPINN loss --------------------------------------------------
def loss(params):
    # PDE residual: -u''(x) - f(x) = 0
    residuals = np.array([
        -d2u_dx2(x, params) - source(x) for x in x_r
    ])
    L_r = np.mean(residuals**2)

    # Boundary conditions: u(0) = u(1) = 0
    L_b = circuit(x_b[0], params)**2 + circuit(x_b[1], params)**2

    return L_r + lambda_bc * L_b

# -- Parameter-shift gradient (exact) ---------------------------
def grad_loss(params, shift=np.pi / 2):
    g = np.zeros(len(params))
    for mu in range(len(params)):
        p_plus  = params.copy(); p_plus[mu]  += shift
        p_minus = params.copy(); p_minus[mu] -= shift
        g[mu] = 0.5 * (loss(p_plus) - loss(p_minus))
    return g
  \end{lstlisting}
\end{frame}

\begin{frame}[fragile]{Training Loop and Results}
  \begin{lstlisting}
# -- Training ----------------------------------------------------
np.random.seed(7)
params = np.random.uniform(0, np.pi, 8)   # 2 layers x 4 params

history = []
# Adam hyperparameters
lr, b1, b2, eps = 0.15, 0.9, 0.999, 1e-8
m, v = np.zeros(8), np.zeros(8)

for step in range(300):
    g = grad_loss(params)
    t = step + 1
    m = b1 * m + (1 - b1) * g
    v = b2 * v + (1 - b2) * g**2
    mc, vc = m / (1 - b1**t), v / (1 - b2**t)
    params -= lr * mc / (np.sqrt(vc) + eps)
    history.append(loss(params))
    if (step + 1) % 75 == 0:
        print(f"Step {step+1:4d}  loss = {history[-1]:.6f}")

# -- Evaluate and compare to exact solution ----------------------
x_test   = np.linspace(0, 1, 100)
u_qpinn  = np.array([circuit(x, params) for x in x_test])
u_exact  = exact(x_test)
l2_error = np.sqrt(np.mean((u_qpinn - u_exact)**2))
print(f"\nL2 error vs. sin(pi*x): {l2_error:.4f}")
  \end{lstlisting}
\end{frame}

\begin{frame}[fragile]{Plotting and Interpretation}
  \begin{lstlisting}
import matplotlib.pyplot as plt

fig, axes = plt.subplots(1, 2, figsize=(11, 3.5))

# -- Solution comparison -----------------------------------------
axes[0].plot(x_test, u_exact, 'k-', lw=2, label=r'Exact $\sin(\pi x)$')
axes[0].plot(x_test, u_qpinn, 'b--', lw=2, label='QPINN (2 qubits, 2 layers)')
axes[0].scatter(x_r, np.zeros_like(x_r), s=20, c='r', zorder=5,
                label='Collocation points')
axes[0].set_xlabel('$x$'); axes[0].set_ylabel('$u(x)$')
axes[0].set_title('QPINN solution: $-u\'\'=\\pi^2\\sin(\\pi x)$')
axes[0].legend(fontsize=8)

# -- Training loss -----------------------------------------------
axes[1].semilogy(history, 'steelblue')
axes[1].set_xlabel('Training step'); axes[1].set_ylabel('Loss')
axes[1].set_title('QPINN training loss (Adam, lr=0.15)')

plt.tight_layout(); plt.show()
  \end{lstlisting}

  \vspace{0.2cm}
  \begin{columns}[T]
    \column{0.48\textwidth}
      \textbf{Expected output:}
      \begin{itemize}
        \item L2 error $\lesssim 0.05$ after 300 steps
        \item Loss decreases by 2--3 orders of magnitude
        \item QPINN traces $\sin(\pi x)$ correctly
      \end{itemize}
    \column{0.48\textwidth}
      \textbf{Improving accuracy:}
      \begin{itemize}
        \item More qubits (wider frequency spectrum)
        \item More variational layers
        \item More collocation points $N_r$
        \item Longer training / lower learning rate
      \end{itemize}
  \end{columns}
\end{frame}

%% ============================================================
\section{Extended Applications}
%% ============================================================

\begin{frame}{Extension: Time-Dependent PDEs}
  \begin{columns}[T]
    \column{0.52\textwidth}
      For $u(x,t)$, encode both space and time:
      \[
        S(x,t) = R_y(\pi x)\otimes R_y(\pi t)
          \otimes\cdots
      \]
      \textbf{1D heat equation:}
      $\partial_t u - \kappa\partial_{xx}u = 0$

      QPINN circuit on $n_x + n_t$ qubits:

      \begin{center}
      \begin{quantikz}[row sep=0.42cm, column sep=0.35cm]
        \lstick{$\ket{0}$\,($x$)}
          & \gate{R_y(\pi x)} & \gate[4]{W(\theta)} & \meter{} \\
        \lstick{$\ket{0}$\,($x$)}
          & \gate{R_y(\pi x)} & \qw                 & \qw \\
        \lstick{$\ket{0}$\,($t$)}
          & \gate{R_y(\pi t)} & \qw                 & \qw \\
        \lstick{$\ket{0}$\,($t$)}
          & \gate{R_y(\pi t)} & \qw                 & \qw
      \end{quantikz}
      \end{center}
    \column{0.44\textwidth}
      \textbf{Input shifts for PDE residual:}
      \begin{align*}
        \partial_t u &\approx
          \tfrac12[u(x,t{+}\tfrac\pi2) - u(x,t{-}\tfrac\pi2)]\\
        \partial_{xx}u &\approx
          u(x{+}\tfrac\pi2,t) - 2u(x,t) + u(x{-}\tfrac\pi2,t)
      \end{align*}
      Each point: 5 circuit evaluations for residual.

      \begin{block}{Initial condition}
        Add IC loss:
        \[
          \mathcal{L}_{\rm IC}
          = \frac1{N_x}\sum_k\bigl|u_\theta(x_k,0)-u_0(x_k)\bigr|^2
        \]
      \end{block}
      \begin{alertblock}{Scaling}
        $d$-dimensional PDE: $n_x + n_t = n$ qubits.
        Classical PINN: exponential in $n$; QPINN: linear.
        Advantage is largest when $d\gg 1$.
      \end{alertblock}
  \end{columns}
\end{frame}

\begin{frame}{The Schrödinger Equation as a QPINN}
  \begin{columns}[T]
    \column{0.52\textwidth}
      The time-dependent Schrödinger equation:
      \[
        i\hbar\,\partial_t\psi(x,t) = H\psi(x,t),\quad
        H = -\frac{\hbar^2}{2m}\partial_{xx} + V(x)
      \]

      Split into real and imaginary parts:
      $\psi = \psi_R + i\psi_I$:
      \begin{align*}
        \partial_t\psi_R &= \phantom{-}\tfrac{\hbar}{2m}\partial_{xx}\psi_I
                           - V\psi_I/\hbar\\
        \partial_t\psi_I &= -\tfrac{\hbar}{2m}\partial_{xx}\psi_R
                           + V\psi_R/\hbar
      \end{align*}
      Two QPINNs: $\psi_R\approx\langle Z_0\rangle_{x,t,\theta_R}$
      and $\psi_I\approx\langle Z_0\rangle_{x,t,\theta_I}$.
    \column{0.44\textwidth}
      \begin{block}{Residual losses}
        \[
          r_R = \partial_t\psi_R
                - \tfrac\hbar{2m}\partial_{xx}\psi_I
                + \tfrac{V}\hbar\psi_I
        \]
        \[
          r_I = \partial_t\psi_I
                + \tfrac\hbar{2m}\partial_{xx}\psi_R
                - \tfrac{V}\hbar\psi_R
        \]
        All derivatives by input-shift;
        $V(x)$ evaluated analytically.
      \end{block}
      \begin{alertblock}{Bridge to Hamiltonian simulation}
        The QPINN Schrödinger formulation is the classical
        simulation of quantum time evolution.
        On a quantum device, Hamiltonian simulation provides
        $e^{-iHt/\hbar}$ exactly --- the QPINN bridges
        continuous PDE solving and digital quantum simulation.
      \end{alertblock}
  \end{columns}
\end{frame}

\begin{frame}{Application to Eigenvalue Problems}
  \begin{columns}[T]
    \column{0.52\textwidth}
      The time-independent Schrödinger equation:
      \[
        H\psi(x) = E\psi(x)
      \]
      is a \textbf{quantum eigenvalue problem}.
      A QPINN approach minimises the Rayleigh quotient:
      \[
        E[\psi_\theta]
        = \frac{\langle\psi_\theta|H|\psi_\theta\rangle}
               {\langle\psi_\theta|\psi_\theta\rangle}
        \ge E_0
      \]
      \textbf{Connection to VQE (week15.tex):}\\
      This is precisely the VQE objective, with the quantum
      ansatz encoding the spatial wavefunction $\psi_\theta(x)$.

      The QFI metric governs the geometry of the wavefunction
      manifold; TDVP (week15.tex) provides the optimal
      projection of the Schrödinger equation.
    \column{0.44\textwidth}
      \begin{block}{Excited states}
        Minimise the augmented loss:
        \[
          \mathcal{L}_n(\theta)
          = \langle H\rangle_\theta
            + \sum_{k<n}\mu_k|\langle\psi_k|\psi_\theta\rangle|^2
        \]
        Penalty $\mu_k$ enforces orthogonality to the
        already-found states $\psi_0,\ldots,\psi_{n-1}$.
        QPINN naturally extends to excited states
        of the Schrödinger, Helmholtz, and
        Sturm-Liouville equations.
      \end{block}
      \begin{alertblock}{QPINN $\leftrightarrow$ VQE}
        VQE on a qubit Hamiltonian (week15.tex) is the
        \emph{discretised} version of a QPINN applied to
        the Schrödinger equation on a spatial grid.
      \end{alertblock}
  \end{columns}
\end{frame}

%% ============================================================
\section{Exercises, Outlook, and References}
%% ============================================================

\begin{frame}{Exercises}
  \begin{enumerate}\itemsep5pt
    \item \textbf{Harmonic oscillator.}
          Modify the Poisson code to solve $u'' + \omega^2 u = 0$
          with $u(0)=0$, $u'(0)=1$.
          Compare to the exact solution $u(x) = \sin(\omega x)/\omega$.

    \item \textbf{Second-order input-shift identity.}
          Prove that for $u_\theta(x) = \langle O\rangle$ with
          $S(x) = e^{-ixP/2}$ (Pauli $P$, eigenvalues $\pm\tfrac12$):
          \[
            \frac{d^2 u_\theta}{dx^2}
            = u_\theta(x+\tfrac\pi2) - 2u_\theta(x) + u_\theta(x-\tfrac\pi2)
          \]

    \item \textbf{Frequency analysis.}
          For a 2-qubit QPINN with angle encoding $R_y(\pi x)$ on each
          qubit and $L=2$ variational layers, list all accessible
          Fourier frequencies $\omega\in\Omega$.
          What is the highest frequency representable?

    \item \textbf{Ritz method.}
          Implement the variational (Ritz) QPINN loss
          $\mathcal{L}_{\rm Ritz} = \sum_k w_k[\tfrac12(u'_\theta)^2 - fu_\theta]$
          for the Poisson problem and compare convergence to the
          strong-form residual loss.

    \item \textbf{Heat equation.}
          Extend the code to solve $\partial_t u = \partial_{xx} u$,
          $x\in(0,1)$, $t\in(0,0.5)$, with $u(x,0)=\sin(\pi x)$
          and $u(0,t)=u(1,t)=0$.
          Exact solution: $u(x,t) = e^{-\pi^2 t}\sin(\pi x)$.
  \end{enumerate}
\end{frame}

\begin{frame}{Outlook: Open Problems and Future Directions}
  \begin{columns}[T]
    \column{0.50\textwidth}
      \begin{block}{Near-term (NISQ)}
        \begin{itemize}
          \item Small QPINNs for 1D/2D PDEs on simulators
          \item Error mitigation for noisy derivatives
          \item Hybrid: QPINN ansatz + classical finite
                element solver
          \item Benchmarking vs.\ classical PINNs on
                spectral-bias-prone problems
        \end{itemize}
      \end{block}
      \vspace{0.2cm}
      \begin{alertblock}{Fundamental trade-off}
        \[
          \underbrace{n \text{ qubits}}_{\text{exp.\ memory}}
          \longleftrightarrow
          \underbrace{2^n \text{ amplitudes}}_{\text{hard to read out}}
        \]
        Read-out bottleneck may eliminate quantum advantage.
      \end{alertblock}
    \column{0.46\textwidth}
      \begin{block}{Long-term directions}
        \begin{itemize}
          \item \textbf{Quantum simulation of PDEs:}
                Hamiltonian simulation provides $e^{-iHt}$
                exactly --- QPINNs as a variational bridge
          \item \textbf{High-$d$ PDEs:} molecular dynamics,
                financial derivatives, kinetic theory
          \item \textbf{Inverse problems:} learn the
                coefficients of a PDE from data
                (combination of QPINN + QBM)
          \item \textbf{Multi-physics:}
                coupled PDE systems via multi-register
                quantum circuits
          \item \textbf{Error correction:}
                fault-tolerant QPINNs for provably
                accurate PDE solvers
        \end{itemize}
      \end{block}
  \end{columns}
\end{frame}

\begin{frame}{Summary: QPINN in Context}
  \begin{center}
  \renewcommand{\arraystretch}{1.3}
  \begin{tabular}{p{2.8cm}p{3.8cm}p{3.8cm}}
    \toprule
    & \textbf{Classical PINN} & \textbf{QPINN} \\
    \midrule
    Ansatz & MLP $u_\theta(x)$ &
             Circuit $\langle O\rangle_{x,\theta}$ \\
    Derivatives & Autograd & Input-shift (exact) \\
    Parameter grad & Backprop & Parameter-shift (exact) \\
    Frequency & Spectral bias & Bounded, explicit $\Omega$ \\
    Memory & $O(N_p)$ & $O(n)$ params, $2^n$ states \\
    BC enforcement & Loss penalty & Same \\
    Optimiser & Adam, L-BFGS & Adam + natural gradient \\
    Geometry & Euclidean & QFI / Fubini-Study \\
    PDE target & Any & Linear PDEs best \\
    Advantage & Production-ready & High-$d$, quantum PDEs \\
    \bottomrule
  \end{tabular}
  \end{center}
  \vspace{0.2cm}
  \begin{alertblock}{The central message}
    A QPINN is a \textbf{physics-informed QNN}: the circuit
    output approximates the PDE solution; all spatial
    derivatives are obtained via the \emph{input-shift rule}
    (a cousin of the parameter-shift rule); training minimises
    the PDE residual using exact quantum gradients.
    The QFI and TDVP from week15.tex provide the natural
    optimisation geometry.
  \end{alertblock}
\end{frame}

\begin{frame}{References}
  \begin{columns}[T]
    \column{0.50\textwidth}
      \textbf{Classical PINNs}
      \begin{enumerate}
        \item M.~Raissi, P.~Perdikaris, G.~E.~Karniadakis,
              \emph{Physics-informed neural networks},
              J.\ Comput.\ Phys.\ \textbf{378}, 686 (2019).
        \item S.~Wang, X.~Yu, P.~Perdikaris,
              \emph{When and why PINNs fail to train},
              J.\ Comput.\ Phys.\ \textbf{449}, 110768 (2022).
      \end{enumerate}
      \textbf{Quantum PINNs}
      \begin{enumerate}\setcounter{enumi}{2}
        \item A.~Kyriienko, A.~E.~Paine, V.~E.~Elfving,
              \emph{Solving nonlinear differential equations
              with differentiable quantum circuits},
              Phys.\ Rev.\ A \textbf{103}, 052416 (2021).
        \item T.~Goto et al.,
              \emph{Universal approximation property of quantum
              machine learning models in quantum-enhanced
              feature spaces},
              Phys.\ Rev.\ Lett.\ \textbf{127}, 090506 (2021).
      \end{enumerate}
    \column{0.46\textwidth}
      \textbf{Fourier analysis of QNNs}
      \begin{enumerate}\setcounter{enumi}{4}
        \item M.~Schuld et al.,
              \emph{Effect of data encoding on the expressive
              power of variational quantum machine learning
              models},
              Phys.\ Rev.\ A \textbf{103}, 032430 (2021).
      \end{enumerate}
      \textbf{Data re-uploading and expressivity}
      \begin{enumerate}\setcounter{enumi}{5}
        \item A.~P\'{e}rez-Salinas et al.,
              \emph{Data re-uploading for a universal quantum
              classifier},
              Quantum \textbf{4}, 226 (2020).
        \item H.-Y.~Huang et al.,
              \emph{Power of data in quantum machine learning},
              Nat.\ Commun.\ \textbf{12}, 2631 (2021).
      \end{enumerate}
      \textbf{QNN foundation (see also week15.tex)}
      \begin{enumerate}\setcounter{enumi}{7}
        \item J.~Stokes et al.,
              \emph{Quantum natural gradient},
              Quantum \textbf{4}, 269 (2020).
        \item M.~Cerezo et al.,
              \emph{Variational quantum algorithms},
              Nat.\ Rev.\ Phys.\ \textbf{3}, 625 (2021).
      \end{enumerate}
  \end{columns}
\end{frame}

\end{document}
